% ============================================================================
\section{Diseño común del programa}\label{sec:programa}
% ============================================================================

Los siete algoritmos están en un único programa en C. Esta sección explica las piezas que
comparten, porque casi todas las decisiones de eficiencia están aquí.

\subsection{Organización en ficheros}

\begin{center}
\begin{tabular}{@{}lp{11.5cm}@{}}
\toprule
Fichero & Contenido \\ \midrule
\cod{scp.h} & Cabecera: define los tipos de datos (instancia, solución, parámetros) y declara
              todas las funciones. No contiene código ejecutable. \\
\cod{base.c} & Herramientas comunes: números aleatorios, reloj, lectura del fichero,
              operaciones sobre una solución, \cod{completar} y \cod{perturbar}. \\
\cod{constructivo.c} & Algoritmo constructivo voraz. \\
\cod{busqueda\_local.c} & Estructuras de entornos $\D1,\D2,\D3$ y búsqueda local (VND). \\
\cod{temple.c} & Temple Simulado. \\
\cod{tabu.c} & Búsqueda Tabú. \\
\cod{grasp.c} & GRASP. \\
\cod{genetico.c} & Algoritmo Genético. \\
\cod{vns.c} & Búsqueda en Entorno Variable. \\
\cod{main.c} & Programa principal: lee los argumentos, llama al algoritmo pedido, comprueba
              la solución y escribe el resultado. \\
\bottomrule
\end{tabular}
\end{center}

\begin{decision}[Decisión: un fichero por algoritmo y un único ejecutable]
Separar cada algoritmo en su fichero permite leerlos de forma independiente. Reunirlos en un
solo ejecutable garantiza que todos usan \emph{exactamente} las mismas rutinas de lectura,
de evaluación y de medida del tiempo, de modo que las comparaciones entre ellos son justas.
\end{decision}

\subsection{Cómo se guarda la instancia}\label{sec:estructuras}

Como solo el 1.5\,\% de la matriz son unos, se guardan dos listas compactas de unos:
\begin{itemize}
  \item \textbf{por filas}: para cada fila $i$, las columnas que la cubren. Sirve para
        responder a ``¿quién puede cubrir esta fila?'';
  \item \textbf{por columnas}: para cada columna $j$, las filas que cubre. Sirve para
        responder a ``¿qué pasa si añado o quito esta columna?''.
\end{itemize}
Cada lista se guarda en dos vectores: uno con todos los índices seguidos y otro que dice
dónde empieza cada fila (o columna). Para el ejemplo de la sección~\ref{sec:ejemplo},
numerando desde 0 como se hace dentro del programa:

\begin{center}
\small
\begin{tabular}{@{}ll@{}}
\toprule
Por filas & Por columnas \\ \midrule
\cod{fila\_col = [0,3,\ 1,3,4,\ 1,2]} & \cod{col\_fil = [0,\ 1,2,\ 2,\ 0,1,\ 1]} \\
\cod{fila\_ini = [0, 2, 5, 7]} & \cod{col\_ini = [0, 1, 3, 4, 6, 7]} \\
\bottomrule
\end{tabular}
\end{center}
Así, las columnas de la fila 1 (la segunda) son \cod{fila\_col[2..4]} $=\{1,3,4\}$, es decir,
las columnas 2, 4 y 5 del enunciado. Este formato se conoce como CSR/CSC y es el habitual
para matrices dispersas.

Además se guarda la matriz completa como una tabla de bytes ($500\times10\,000=5$~MB), que
permite preguntar al instante ``¿cubre la columna $j$ la fila $i$?'' (macro \cod{CUBRE}). Solo
se usa para esa consulta puntual, nunca se recorre entera.

\fragmento{scp.h}{52}{66}

\subsection{Cómo se guarda una solución}\label{sec:solucion}

Una solución no es solo el conjunto de columnas. Se mantienen \emph{siempre actualizados}:
\begin{itemize}
  \item \cod{en[j]}: 1 si la columna $j$ está elegida;
  \item \cod{lista} y \cod{pos}: las columnas elegidas, una tras otra, y la posición de cada
        una en esa lista (permite recorrer solo las elegidas y borrar una en tiempo constante);
  \item \cod{cub[i]}: \textbf{cuántas columnas elegidas cubren la fila $i$};
  \item \cod{nsin}: cuántas filas están sin cubrir (la solución es factible si y solo si
        \cod{nsin} $=0$);
  \item \cod{z}: el coste total.
\end{itemize}

\begin{idea}[Idea clave: evaluación incremental]
Gracias al contador \cod{cub}, añadir o quitar una columna solo obliga a tocar las filas de
\emph{esa} columna (unas 7 u 8 de media), no las 500 filas ni los 75\,000 unos. Como las
metaheurísticas hacen millones de movimientos, esta es la decisión que más influye en la
velocidad del programa. También da gratis las dos preguntas que se hacen continuamente:
\emph{¿es factible?} (\cod{nsin == 0}) y \emph{¿es redundante la columna $j$?} (todas sus
filas tienen \cod{cub} $\ge 2$).
\end{idea}

\fragmento{base.c}{313}{328}

\subsection{Todas las soluciones son siempre factibles}

\begin{decision}[Decisión: no trabajar con soluciones infactibles]
Hay dos formas de tratar la restricción ``todas las filas cubiertas'': (a) permitir soluciones
que dejen filas sin cubrir y penalizarlas en la función objetivo, o (b) no salir nunca del
conjunto factible. Se ha elegido (b) en los cinco algoritmos. Razones: no hay que inventar ni
calibrar un peso de penalización; cualquier solución que el algoritmo tenga entre manos se
puede devolver tal cual; y en este problema recuperar la factibilidad es muy fácil y barato
(basta añadir columnas). Cuando un operador rompe la factibilidad (al quitar columnas, al
cruzar dos soluciones) se \textbf{repara} inmediatamente.
\end{decision}

La reparación y la construcción son, en el fondo, la misma operación: \emph{cubrir las filas
que estén sin cubrir}. Por eso existe una única función, \cod{completar}, que usan el
constructivo (partiendo de la solución vacía), GRASP, el temple, el genético y la VNS
(partiendo de una solución a la que le faltan columnas). Se explica en la
sección~\ref{sec:constructivo}.

\subsection{La perturbación \texorpdfstring{$\Om{k}$}{Omega\_k}}\label{sec:perturbar}

Varios algoritmos necesitan ``una solución parecida a $S$ elegida al azar''. Se obtiene con la
función \cod{perturbar}:
\begin{enumerate}
  \item se quitan $k$ columnas de $S$ elegidas al azar (quedan filas sin cubrir);
  \item se repara con \cod{completar}, prohibiendo que vuelvan a entrar justo las $k$ que se
        acaban de quitar;
  \item se eliminan las columnas que hayan quedado redundantes.
\end{enumerate}
Llamaremos $\Om{k}(S)$ al conjunto de soluciones que pueden salir así. Con $k$ pequeño la
nueva solución se parece mucho a $S$; con $k$ grande, poco.

\begin{decision}[Decisión: vetar las columnas recién quitadas]
Sin el veto del paso 2, el voraz volvería a meter muchas veces las mismas columnas que se
acaban de sacar (por algo estaban en la solución) y la ``vecina'' sería la propia $S$. El
veto solo dura lo que dura la reparación.
\end{decision}

\fragmento{base.c}{577}{595}

\subsection{Números aleatorios, tiempo y precisión}

\begin{itemize}
  \item \textbf{Generador aleatorio propio.} No se usa la función \cod{rand()} de C porque da
        secuencias distintas en cada sistema operativo. Se usa un generador \emph{xorshift}
        de 64 bits escrito en \cod{base.c}: con la misma semilla, el programa hace
        exactamente lo mismo en cualquier ordenador. Eso hace los experimentos
        \textbf{reproducibles} y permite depurar.
  \item \textbf{Criterio de parada por tiempo de CPU.} Todas las metaheurísticas se detienen
        al consumir un tiempo de procesador fijado (5 segundos por defecto), medido con
        \cod{clock()}. Es el criterio más justo para comparar algoritmos cuyas iteraciones
        cuestan cantidades muy distintas (una iteración del temple dura microsegundos; una
        generación del genético, milisegundos). Se mide tiempo de CPU y no de reloj para que
        no influya lo ocupado que esté el ordenador.
  \item \textbf{Tolerancia numérica.} Los costes son números reales. Para que un error de
        redondeo de $10^{-13}$ no se tome por una mejora, ``$a$ mejora a $b$'' significa
        $a<b-10^{-9}$ (constante \cod{EPS}). Además, el coste de la solución final se vuelve a
        sumar desde cero antes de escribirlo.
  \item \textbf{Comprobación independiente.} Antes de escribir el resultado,
        \cod{sol\_comprobar} recuenta desde cero la cobertura de cada fila y el coste, sin
        fiarse de los contadores. El programa solo declara una solución ``factible'' si pasa
        esa comprobación. En las 1548%
 ejecuciones de esta memoria
        la pasan todas.
\end{itemize}
